\documentclass[11pt]{article}
\usepackage[utf8]{inputenc}
\usepackage[T1]{fontenc}
\usepackage{amsmath,amssymb}
\usepackage{geometry}
\geometry{a4paper, margin=2cm}

\begin{document}

\begin{center}
    {\LARGE \textbf{Gabarito Recuperação - Matemática Básica}}\\[6pt]
    {\large Resoluções detalhadas}
\end{center}

\bigskip

\begin{enumerate}

\item Um estudante registrou quantas horas dormiu em 5 dias: $6{,}5;\;7{,}0;\;6{,}8;\;7{,}2;\;7{,}0$ horas.  
Calcule a média de horas dormidas. Mostre como fez o cálculo.

\textbf{Cálculo:}
\[
\text{Soma} = 6{,}5 + 7{,}0 + 6{,}8 + 7{,}2 + 7{,}0 = 34{,}5
\]
\[
\text{Média} = \frac{34{,}5}{5} = 6{,}9
\]

\textbf{Resposta:} A média é \(\mathbf{6{,}9}\) horas.

\bigskip

\item As notas de uma aluna em cinco atividades foram: $8{,}0;\;7{,}5;\;8{,}5;\;7{,}0;\;8{,}0$.  
Determine a mediana dessas notas. Explique os passos.

\textbf{Cálculo:} ordenar os valores:
\[
7{,}0,\;7{,}5,\;8{,}0,\;8{,}0,\;8{,}5.
\]
Com 5 observações a mediana é o termo do meio (3.º): \(8{,}0\).

\textbf{Resposta:} Mediana = \(\mathbf{8{,}0}\).

\bigskip

\item As idades de seis crianças em uma creche são: $4,\,5,\,4,\,6,\,4,\,5$ anos.  
Qual é a moda dessas idades? Mostre sua justificativa.

\textbf{Cálculo:} contagem das frequências:
\[
4 \to 3\ \text{vezes},\quad 5 \to 2,\quad 6 \to 1.
\]
Moda = valor com maior frequência = \(4\).

\textbf{Resposta:} Moda = \(\mathbf{4}\) anos.

\bigskip

\item Um corredor registrou seus tempos (em minutos) em quatro treinos: $12,\;11,\;13,\;12$.  
Calcule a média dos tempos.

\textbf{Cálculo:}
\[
\text{Soma} = 12 + 11 + 13 + 12 = 48
\]
\[
\text{Média} = \frac{48}{4} = 12
\]

\textbf{Resposta:} Média = \(\mathbf{12}\) minutos.

\bigskip

\item O peso (em kg) de cinco pacotes de frutas foi: $1{,}8;\;2{,}0;\;1{,}9;\;2{,}1;\;2{,}0$.  
Determine a mediana dos pesos. Explique.

\textbf{Cálculo:} ordenar:
\[
1{,}8,\;1{,}9,\;2{,}0,\;2{,}0,\;2{,}1.
\]
Com 5 observações a mediana é o termo do meio (3.º): \(2{,}0\).

\textbf{Resposta:} Mediana = \(\mathbf{2{,}0}\) kg.

\bigskip

\item A média das idades de 5 pessoas é $25$ anos. Quatro delas têm $22,\;24,\;26,\;28$ anos.  
Qual é a idade da quinta pessoa? Mostre o raciocínio.

\textbf{Cálculo:} soma total necessária:
\[
5 \times 25 = 125.
\]
Soma das quatro dadas:
\[
22 + 24 + 26 + 28 = 100.
\]
Logo a quinta idade é
\[
x = 125 - 100 = 25.
\]

\textbf{Resposta:} A quinta pessoa tem \(\mathbf{25}\) anos.

\bigskip

\item As notas de cinco estudantes foram: $6{,}5;\;7{,}0;\;8{,}0;\;7{,}5;\;x$.  
Sabendo que a mediana é $7{,}0$, descubra o valor de $x$. Explique o motivo.

\textbf{Cálculo e argumento:} com 5 valores a mediana é o 3.º menor. Ordenando os quatro conhecidos temos: \(6{,}5,\;7{,}0,\;7{,}5,\;8{,}0\). Para que o elemento do meio (3.º) seja \(7{,}0\), após inserir \(x\) na ordenação o terceiro valor deve continuar sendo \(7{,}0\). Isso ocorre sempre que \(x \le 7{,}0\) (por exemplo \(x=7{,}0\) ou qualquer valor menor que \(7{,}0\) mas $\ge 6{,}5$ ou menor que 6{,}5 — ainda assim o 3.º continuará sendo 7{,}0).  

\textbf{Conclusão:} qualquer \(x\) tal que \(\mathbf{x \le 7{,}0}\) mantém a mediana igual a \(7{,}0\). Em particular, \(x=7{,}0\) é uma solução válida.

\bigskip

\item Em uma pesquisa sobre o número de animais de estimação por família, os dados foram: $1,\,2,\,1,\,3,\,x,\,2$.  
Sabendo que a moda é $1$, determine o valor de $x$. Justifique.

\textbf{Cálculo:} frequências atuais: \(1\to2\), \(2\to2\), \(3\to1\). Para que a moda seja exclusivamente \(1\) (isto é, \(1\) seja o valor mais frequente), precisamos que a contagem de \(1\) aumente para mais que 2; assim devemos ter \(x=1\), passando a frequência de \(1\) para 3.

\textbf{Resposta:} \(\mathbf{x = 1}\).

\bigskip

\item O tempo (em minutos) gasto por quatro funcionários para concluir uma tarefa foi: $45,\;50,\;55,\;x$.  
Se a média foi de $52$ minutos, determine o valor de $x$. Mostre os cálculos.

\textbf{Cálculo:} soma total necessária:
\[
4 \times 52 = 208.
\]
Soma conhecida:
\[
45 + 50 + 55 = 150.
\]
Logo
\[
x = 208 - 150 = 58.
\]

\textbf{Resposta:} \(\mathbf{x = 58}\) minutos.

\bigskip

\item As notas obtidas por sete alunos em uma prova foram: $6{,}0;\;6{,}5;\;7{,}0;\;7{,}5;\;8{,}0;\;8{,}5;\;x$.  
Sabendo que a mediana é $7{,}5$, determine o valor de $x$. Explique seu raciocínio.

\textbf{Cálculo e argumento:} com 7 valores a mediana é o 4.º menor. Na lista já ordenada (sem \(x\)) o 4.º elemento é \(7{,}5\). Para que a mediana continue \(7{,}5\) após inserir \(x\), \(x\) deve ser maior ou igual a \(7{,}5\) (se \(x<7{,}5\) ele empurraria \(7{,}5\) para a 5.ª posição). Portanto qualquer \(x\) com \(x \ge 7{,}5\) mantém a mediana em \(7{,}5\).  

\textbf{Conclusão:} \(x\) deve satisfazer \(\mathbf{x \ge 7{,}5}\). Em particular \(x=7{,}5\) é uma solução simples.

\end{enumerate}

\end{document}

